% Part: lambda-calculus
% Chapter: church-rosser
% Section: definitions-and-properties

\documentclass[../../../include/open-logic-section]{subfiles}

\begin{document}

\olfileid[fr]{lam}{cr}{dap}
\olsection{Définition et propriétés}

Dans ce chapitre, nous introduisons la
propriété de Church--Rosser et quelques-unes
de ses conséquences.

\begin{defn}[Propriété de Church--Rosser, CR]
Une relation $\xredone$ sur les termes
vérifie la \emph{propriété de Church--Rosser}
si, chaque fois que $M \xredone P$ et
$M \xredone Q$, il existe un terme~$N$ tel
que $P \xredone N$ et $Q \xredone N$.
\end{defn}

On peut voir le lambda-calcul comme un modèle
de calcul dans lequel les termes en forme
normale sont des « valeurs » et une relation
de réduction fournit les « règles de calcul ».
La propriété de Church--Rosser assure que
deux calculs issus du même terme peuvent
encore se rejoindre. S'ils atteignent une
forme normale, celle-ci est donc unique.

Prenons un exemple d'algèbre élémentaire :
on peut calculer $4 \times (1+2) + 3$ de
plusieurs manières. Si l'on commence par
réduire $1+2$ à~$3$, on obtient
$4 \times 3+3$. Si l'on distribue d'abord
$4 \times (1+2)$, on obtient
$4 \times 1+4 \times 2+3$. Les deux
expressions se réduisent ensuite à~$12+3$.

Si $\xredone$ est la réduction~$\beta$,
la propriété de Church--Rosser entraîne
qu'un terme qui possède une forme normale
n'en possède qu'une. En effet, supposons
que $M$ se réduise à deux formes normales
$P$ et~$Q$. Il existe alors un terme~$N$
auquel $P$ et $Q$ se réduisent tous deux.
Puisque $P$ et $Q$ sont en forme normale,
ces réductions ne peuvent être que réflexives :
$P$, $Q$ et $N$ sont identiques. On peut
donc parler de \emph{la} forme normale d'un terme.

Si l'on considère le lambda-calcul comme un
modèle de calcul, la forme normale représente
le « résultat final » d'un calcul qui part
du terme. Le corollaire précédent dit qu'il
y a au plus un tel résultat, comme le calcul
de $4 \times (1+2)+3$ donne~$15$.

\begin{thm}\ollabel{thm:str}
Si une relation $\xredone$ vérifie la
propriété de Church--Rosser et si $\xred$
est la plus petite relation transitive qui
contient $\xredone$, alors $\xred$ vérifie
elle aussi la propriété de Church--Rosser.
\end{thm}

\begin{proof}
Supposons
\begin{align*}
  M & \xredone P_1 \xredone \dots \xredone P_m
    \quad\text{et}\\
  M & \xredone Q_1 \xredone \dots \xredone Q_n.
\end{align*}
Construisons une grille~$N$ de termes de
hauteur $m+1$ et de largeur $n+1$.
Le terme $N_{i,j}$ occupe la ligne~$i$
et la colonne~$j$. Nous voulons que
$N_{i,j}\xredone N_{i+1,j}$ et
$N_{i,j}\xredone N_{i,j+1}$ partout où
ces cases existent. Posons :
\begin{align*}
  N_{0,0} &= M \\
  N_{i,0} &= P_i && \text{si }1\le i\le m \\
  N_{0,j} &= Q_j && \text{si }1\le j\le n \\
  \intertext{et, pour les cases intérieures :}
  N_{i,j} &= R.
\end{align*}
Ici, $R$ est un terme tel que
$N_{i-1,j}\xredone R$ et
$N_{i,j-1}\xredone R$. Pour chaque case
intérieure, les deux termes précédents
sont issus en un pas de $N_{i-1,j-1}$;
la propriété de Church--Rosser assure
donc l'existence de~$R$. On construit
ainsi la grille de proche en proche.

Enfin, $N_{m,0}\xredone\dots\xredone N_{m,n}$
et $N_{0,n}\xredone\dots\xredone N_{m,n}$.
En posant $P=P_m=N_{m,0}$ et
$Q=Q_n=N_{0,n}$, ces deux chaînes montrent
que $P\xred N_{m,n}$ et $Q\xred N_{m,n}$.
La propriété de Church--Rosser de $\xred$
en découle.
\end{proof}

\end{document}
